\chapter{When Things Go Wrong}\label{ch-errors}

A correct program can still meet things that go wrong: a line that
should hold a number holds a word, a divisor turns out to be zero, a
file is not there. Zeno calls these \emph{errors}, and gives a program
three words to deal with them: \verb|raise| to say that something has
gone wrong, \verb|try| and \verb|except| to deal with it, and
\verb|raises| to say in a function's heading what it may raise.

\section{Bugs and errors}

Zeno draws a line between two kinds of thing that go wrong.

A \emph{bug} is a mistake in the program itself: a misspelt name,
\verb|"abc" - 1|, calling something that is not a function. The
checker finds almost all of them before the program runs, and the
program does not start until they are fixed.

An \emph{error} is something that can go wrong in a correct program:
bad input, a division by a value that turned out to be zero, a
position that came from outside. These happen while the program runs,
and the program can catch them and carry on.

\begin{why}
Catching a bug hides it. A program that catches "this name does not
exist" and carries on has a mistake in it that nobody will ever see.
Since the checker reports bugs before the program runs, there is
nothing of that kind left to catch. (Decision 4.1.)
\end{why}

\section{An error that is not caught}

An error that nothing catches stops the program, with its kind, its
message and the line it came from:

\begin{zeno}[fails=yes]
let readings = [20.5, 21.0]
print(readings[2])
\end{zeno}

Here \verb|IndexError| is the \emph{kind} of error, and the rest is
its \emph{message}. Every error has both.

\section{try and except}

Code that may raise an error goes in a \verb|try| block. If an error
is raised inside it, the rest of the block is skipped, and the
\verb|except| that names its kind runs instead:

\begin{zeno}
for text in ["40", "fast", "7"] {
    try {
        let speed = int(text)
        print("speed", speed)
    } except ValueError {
        print("not a number:", text)
    }
}
\end{zeno}

If nothing is raised, the \verb|except| block does not run at all.

\verb|except Kind as e| gives the error a name. \verb|e.message| is
its message and \verb|e.kind| the name of its kind:

\begin{zeno}
try {
    print(int("fast"))
} except ValueError as e {
    print(e.kind, "says:", e.message)
}
\end{zeno}

A \verb|try| can have several \verb|except|s. The first that fits is
the one that runs:

\begin{zeno}
function show(xs: list[int], at: int, by: int) {
    try {
        print(xs[at] // by)
    } except IndexError as e {
        print("position:", e.message)
    } except ZeroDivisionError {
        print("cannot divide by zero")
    }
}

show([10, 20], 1, 5)
show([10, 20], 9, 5)
show([10, 20], 0, 0)
\end{zeno}

An error that no \verb|except| names goes on as if the \verb|try| were
not there, up to whoever called this code, and so on until something
catches it or the program stops.

\section{The kinds of error}

The kinds that Zeno itself raises are these:

\begin{center}
\begin{tabular}{lp{9cm}}
\toprule
Kind & Raised when \\
\midrule
\verb|ValueError| & a value is the right type but cannot be used:
  \verb|int("fast")|, \verb|2 ** -1| \\
\verb|IndexError| & a position is outside a list or a text \\
\verb|KeyError| & a dict has no such key \\
\verb|ZeroDivisionError| & a number is divided by zero \\
\verb|OverflowError| & a whole number is too big for its type \\
\verb|FileError| & a file cannot be read or written \\
\verb|Error| & anything else; and \verb|except Error| catches every kind \\
\bottomrule
\end{tabular}
\end{center}

\verb|except Error| catches any error at all, which is useful at the
very top of a program, to report what went wrong and stop tidily:

\begin{zeno}
let big = 9223372036854775807
try {
    print(big + 1)
} except Error as e {
    print("stopped:", e.kind, e.message)
}
\end{zeno}

\hyperref[app-errors]{Appendix F} lists every message that the
library and the runtime give.

\section{raise}

A program raises an error of its own with \verb|raise| and the kind,
with a message that says what went wrong:

\begin{zeno}
function check_speed(speed: int) -> int raises ValueError {
    if speed < 0 or speed > 100 {
        raise ValueError(f"a speed is 0 to 100, and {speed} is not")
    }
    return speed
}

for s in [40, 250] {
    try {
        print("set to", check_speed(s))
    } except ValueError as e {
        print(e.message)
    }
}
\end{zeno}

A message is for the person who reads it when something goes wrong.
The best ones say what was expected and what was found, as Zeno's own
messages do.

Only an error can be raised, not a piece of text:

\begin{zeno}[fails=yes]
raise "something went wrong"
\end{zeno}

\subsection{Raising again}

An \verb|except| block that cannot fully deal with an error can do
what it can, then raise the same error again, so that the caller sees
it:

\begin{zeno}[fails=yes]
try {
    print(int("fast"))
} except ValueError as e {
    print("noted:", e.message)
    raise e
}
\end{zeno}

\section{raises: what a function may raise}

\verb|raises| after a function's heading lists the errors it may
raise, so that whoever calls it can see them without reading the body:

\begin{zeno}
function read_volts(text: str) -> float raises ValueError {
    return float(text)
}
print(read_volts("3.3"))
\end{zeno}

It follows the rule of every type in Zeno: it can be left out, and
then Zeno works it out from the body and passes it on to callers. When
it is written, it must be complete. A function that says what it
raises and then leaves something out is told so:

\begin{zeno}[fails=yes]
function read_volts(text: str) -> float raises KeyError {
    return float(text)
}
\end{zeno}

In strict mode (\verb|--strict|, \hyperref[ch-types]{Types, Written
and Worked Out}), as in every compiled program, it must always be
written, so that every heading is a complete promise:

\begin{zeno}[fails=yes, cmd=zeno check --strict {file}]
function read_volts(text: str) -> float {
    return float(text)
}
print(read_volts("3.3"))
\end{zeno}

\subsection{What counts}

Only errors raised \emph{on purpose} must be listed: a \verb|raise|, a
call to a function that raises, and \verb|int()| or \verb|float()| on
text, which fail on bad input. The errors that operators raise, such
as dividing by zero, a position out of range or an overflow, can be
caught, but need not be listed.

\begin{why}
A division or an index can fail on almost any line. If those had to be
listed, nearly every function would say \verb|raises ZeroDivisionError,
IndexError|, and a label that is on everything tells nobody anything.
Java's checked exceptions are widely disliked for exactly that.
(Decision 1.17.)
\end{why}

\section{An error that carries more}

An error often needs to carry more than a message: how long was waited,
which pin, what code a device sent back. An object of any class can be
raised, with whatever fields it needs:

\begin{zeno}
class Timeout {
    after_ms: int
    device: str
}

function wait_for(device: str) raises Timeout {
    raise Timeout(30, device)
}

try {
    wait_for("sensor")
} except Timeout as t {
    print("gave up on", t.device, "after", t.after_ms, "ms")
}
\end{zeno}

\verb|except Error as e| catches such an object too. Then
\verb|e.kind| is the class's name, and \verb|e.message| is the object
as text:

\begin{zeno}
class Timeout {
    after_ms: int
}

try {
    raise Timeout(30)
} except Error as e {
    print(e.kind, "/", e.message)
}
\end{zeno}

\begin{why}
A class already means "these named values, together", so there is no
second way to declare an error. (Decision 1.17.)
\end{why}

\section{What is not here}

Zeno has no \verb|finally|. Code that must run whether or not an error
was raised is written after the \verb|try|, with an \verb|except| that
catches what must be caught.

On a microcontroller, a compiled program can carry its errors as plain
values instead, with no jumps at all, as firmware coding standards
require. Nothing in the program changes for that; it is a choice made
when building, and \hyperref[ch-firmware]{Firmware} explains it.

\begin{why}
They are not C++ exceptions: there are no unwinding tables, and code
that does not use \verb|try| pays nothing for it. (Decisions 1.17 and
5.5i.)
\end{why}

\begin{summary}
A bug is found before the program runs; an error happens while it
runs. \verb|raise Kind("message")| raises an error, and
\verb|try { } except Kind as e { }| catches it, with \verb|e.kind| and
\verb|e.message|. \verb|except Error| catches every kind. An object of
any class can be raised. \verb|raises| lists what a function raises on
purpose: worked out when left out, complete when written, and always
written in strict mode.
\end{summary}

\begin{tryit}
\begin{enumerate}
\item Write a loop that asks for numbers with \verb|input()| until it
  is given one that \verb|int| accepts.
\item Write \verb|function percent(part: int, whole: int) -> float|
  that raises a \verb|ValueError| with a clear message when
  \verb|whole| is 0, and catch it.
\item Make a class \verb|SensorFault| with a pin and a code, raise it
  from a function, and catch it once by its own name and once as
  \verb|Error|.
\item Remove \verb|raises ValueError| from \verb|check_speed| and run
  \verb|zeno check --strict|. Then write \verb|raises KeyError|
  instead. What does each message ask for?
\end{enumerate}
\end{tryit}
